% 2026-11-10 / v22
\section{Artin基底における線束生成残余圏と行列因子分解}
\label{am:sec:comparison}

本節では絶対特異圏を用いず、周囲の代数上perfectな複体だけを用いて、
有限階数graded MFと線束から生成される射影残余圏を比較する。
第\ref{cg:sec:comparison}節および第\ref{fo:sec:formal-cubic}節のMF同値とは別の比較である。
非正則Artin基底上で通常の有界coherent圏全体へOrlovの定理を適用することはしない。
相対的なstabilizationの加群論的入力は
Chen\cite[Theorem 2.10, Corollary 2.11, Theorem 7.2]{AMChen}である。
同論文の定理は三角圏についてのものであるため、以下ではexact圏とdg商を通じて
実際のenhanced関手を構成する。
本節の組合せと同じ一般性での先行研究上の初出は未確認である。

\subsection{入力と相対的な特異圏}

変数の数を$d\ge5$、Artin切断の次数を$N\ge0$とし、
\[
 R_N=\CC[h]/(h^{N+1}),\qquad k=R_N/(h),\qquad
 a=d-3,\qquad i=1-a=4-d
\]
と置く。$A=\bigoplus_{j\ge0}A_j$は中央$R_N$-代数で、$A_0=R_N$、
各$A_j$は有限自由$R_N$-加群、
$\overline A=A\otimes_{R_N}k\simeq k[z_1,\ldots,z_d]$と仮定する。
$\Omega\in Z(A)_3$の剰余$F$は滑らかな射影cubicを定め、$B=A/(\Omega)$とする。
以下の構成は内部次数を保つ。本節のMFと厚いdg部分圏は冪等完備化して扱う。

この入力から、左右noetherian性、$R_N$の左右$A$-加群としての有限自由Koszul型分解、
$\Omega$の正則性、$B_j$の自由性が従う。
これらは正則基底上の有限global dimensionと区別する。
実際、$h$による有限filtrationの随伴次数環がnoetherianなので$A,B$はnoetherianである。
補題\ref{lp:lem:resolution}の証明は基底の正則性を使わず、
各次数の自由性とNakayamaの補題だけを用いるので、$R_N$上にも適用できる。
多項式環のKoszul分解を持ち上げると
\begin{equation}\label{am:eq:relative-koszul}
 P_j=A(-j)^{\binom dj}\quad(0\le j\le d)
\end{equation}
という分解を得る。各内部次数で$R_N$-分裂し、全てのArtin商への底変更後も完全である。
$\Omega:A_{j-3}\to A_j$は剰余上で単射なので、有限自由加群の行列の単元小行列式を用いて
$R_N$上でも分裂単射となる。従って$B_j$は自由である。
この議論から$\operatorname{gldim}A<\infty$は結論しない。

\begin{definition}[周囲でperfectな特異圏]\label{am:def:relative-singularity}
標準dg enhancement内で
\[
 \mathcal D_{B/A}
 =\{M\in D^b(\operatorname{gr}B):
       \operatorname{Res}_A M\in\operatorname{Perf}(\operatorname{gr}A)\}
\]
を取り、
\[
 \mathcal S_{B/A}
 =\bigl(\mathcal D_{B/A}/\operatorname{Perf}(\operatorname{gr}B)\bigr)^{\mathrm{idem}}
\]
をそのdg商の冪等完備化とする。
$\operatorname{Perf}(\operatorname{gr}B)\subset\mathcal D_{B/A}$は
$0\to A(-3)\xrightarrow{\Omega}A\to B\to0$から従う。
\end{definition}

ここでのperfect性は周囲の\emph{$A$上}のperfect性である。
$B(j)$やその内部切断は、$R_N$上では一般に無限階数であり、
$\operatorname{Perf}(R_N)$の対象だとは述べない。
また$\mathcal S_{B/A}$を絶対特異圏$D_{\rm sg}(B)$と同一視しない。

\begin{lemma}[相対dg stabilization]\label{am:lem:stabilization}
上の入力に対し、内部shiftと三次周期性を保つdg Morita同値
\begin{equation}\label{am:eq:stabilization}
 \operatorname{cok}:
 \MF^{\ZZ}(A,\Omega)\xrightarrow{\sim}\mathcal S_{B/A}
\end{equation}
が存在する。そのMorita逆を$\operatorname{stab}$と書く。
この同値は有限自由MFの余核とその周期自由分解から誘導される。
Artin商への底変更に関して、これらの構成は自然である。
\end{lemma}

\begin{proof}
まず、次数付き有限射影$A$-加群は次数付き有限自由である。
実際、そのaugmentation商は局所環$R_N$上の次数付き有限射影加群なので自由であり、
その基底を持ち上げた自由加群からの写像に次数付きNakayamaを適用すればよい。

$\mathcal E$を、有限生成次数付きGorenstein-projective $B$-加群のうち
$A$上有限射影次元を持つもののexact圏とする。
Chenのchange-of-rings定理により、非零のこのような加群は
$A$上の射影次元が一である。
正則中心元の場合の同論文の証明は、全ての射と分解を斉次に取れば次数付き圏でも成立する。
特に$\mathcal E$はFrobenius exact圏であり、そのprojective-injective対象は
有限次数付き射影$B$-加群である。
syzygyとcosyzygyで有限$A$-射影次元が保たれること、およびその次元が再び一になることが、
このFrobenius部分圏への制限を保証する。

$\mathcal F$を有限自由graded MFのexact圏とする。
exact列は二つの成分で$A$-分裂する列であり、projective-injective対象は
$(1,\Omega)$、$(\Omega,1)$型の可縮factorizationの直和因子である。
余核はexact関手
\[
 \operatorname{cok}:\mathcal F\longrightarrow\mathcal E
\]
を与える。余核加群の$A$-自由分解の長さは一である。
逆に$M\in\mathcal E$の斉次分解
$0\to P_1\xrightarrow{\varphi}P_0\to M\to0$を取ると、
$\Omega M=0$より$\Omega\Id_{P_0}$は$\varphi$を通って一意にfactorする。
$\varphi$の単射性から他方の積も$\Omega\Id_{P_1}$となり、MFを得る。
射の持上げとhomotopyの同定を含め、これはChenのTheorem 7.2の余核同値である。
従って余核関手はstable categoryの三角同値を与える。

この三角同値だけでdg同値を宣言せず、次のdg商を取る。
exact関手から実際のdg関手
\begin{equation}\label{am:eq:exact-dg-quotients}
 \frac{D^b_{\rm dg}(\mathcal F)}{K^b_{\rm dg}(\operatorname{proj\!-\!inj}\mathcal F)}
 \longrightarrow
 \frac{D^b_{\rm dg}(\mathcal E)}{K^b_{\rm dg}(\operatorname{proj}B)}
\end{equation}
を得る。Frobenius exact圏の右辺・左辺の$H^0$は対応するstable categoryである。
これはprojective-injective分解を用い、有限複体をその高いsyzygyへ移すことで確認できる。
従って\eqref{am:eq:exact-dg-quotients}は$H^0$で同値であり、
shiftを含む任意の対象対へ同じ事実を適用すると全Hom cohomologyで同型となる。

左側のdg商から元のMF dg圏へは、有限複体の有限totalizationを取る。
成分別split exactな複体のtotalizationは可縮である。
具体的に、横微分$d$と横contraction$s$を
$ds+sd=1$、$s^2=0$とし、符号を含む縦MF微分を$\delta$とする。
$d\delta+\delta d=0$、$\delta^2=w$で、$w$は中央曲率である。
$s$は横次数を下げるので、横方向の有限性から
$1+\delta s$、$1+s\delta$は有限和で逆を持つ。
\[
 q=s(1+\delta s)^{-1}=(1+s\delta)^{-1}s
\]
と置けば、直接展開により
\[
 (1+s\delta)\bigl(1-(d+\delta)q-q(d+\delta)\bigr)
 (1+\delta s)=-sws=0
\]
となる。最後は中央性と$s^2=0$を用いた。
従って$(d+\delta)q+q(d+\delta)=1$であり、
totalizationは実際に可縮である。この計算は曲率の冪零性を必要としない。
可縮MFも零へ送られるため、totalizationは左側のdg商を通る。
$H^0$では通常のMF homotopy categoryへの同値であり、同様に全Homで準同型となる。

右側から$\mathcal S_{B/A}$へのdg関手は、$\mathcal E$の$B$-加群圏への包含から得る。
完全射影分解を用いると、Gorenstein-projective加群間のstable Homは特異圏のHomに一致する。
同じ分解のshiftを用いることで全Hom cohomologyについて一致する。
本質的全射性も具体的に確認できる。
$M\in\mathcal D_{B/A}$の有界上の有限生成$B$-射影分解を十分低い位置で切断すると、
特異圏で$M$はある$B$-加群$L$のshiftに同型となり、$L$は$A$上有限射影次元を持つ。
ChenのTheorem 2.10より$L$は$B$上有限Gorenstein-projective次元を持つので、
さらに高い$B$-syzygyを取れば$\mathcal E$に属する。
これで右側のdg商と$\mathcal S_{B/A}$の同値を得る。
三つの実際のdg関手のzigzagにより\eqref{am:eq:stabilization}が得られる。

最後に底変更を確認する。
MFの余核$M$の各内部次数の分解は有限自由$R_N$-加群間の列である。
剰余上で最初の射は単射である。実際、二つのfactorization mapの積は
非零多項式$F$の乗法なので、どちらのmapも単射である。
単元小行列式により各次数の余核$M_j$は$R_N$-自由となる。
従って余核、exact列、有限totalizationはArtin商への底変更と両立する。
周期分解とdg商から得た比較は、Morita localizationの普遍性を通じて自然となる。
選択したcontractionそのものが全ての底変更と厳密に可換することは必要としない。
内部shiftと周期性は全ての段階で保たれ、Morita逆もこれらと両立するように取れる。
\end{proof}

\subsection{Artin基底上の線束Homと有限window}

$\pi:\operatorname{Gr}B\to\operatorname{QGr}B$を射影商とし、
$\mathcal O(j)=\pi B(j)$と置く。
導来sectionを$\mathbf R\omega$と書く。
この段階では$D(\operatorname{QGr}B)$内で線束から生成するdg部分圏を用い、
全ての有界coherent複体についての底変更定理は要求しない。

\begin{lemma}[線束のsectionと底変更]\label{am:lem:line-bundles}
自然なderived Homの底変更は各Artin商と両立し、
\begin{equation}\label{am:eq:line-hom}
 \operatorname{Ext}^q(\mathcal O(r),\mathcal O(s))=
 \begin{cases}
 B_{s-r}&q=0,\\
 B_{r-s-a}^{\vee}&q=d-2,\\
 0&\text{その他}
 \end{cases}
\end{equation}
である。ここで$(-)^\vee=\operatorname{Hom}_{R_N}(-,R_N)$、負添字は零とする。
また$0\le j\le d-1$について自然に
\begin{equation}\label{am:eq:window-sections}
 (\mathbf R\omega\mathcal O(j))_{\ge i}\simeq B(j)_{\ge i}
\end{equation}
となり、この同定もArtin商への底変更と両立する。
\end{lemma}

\begin{proof}
式\eqref{am:eq:relative-koszul}を$A$で双対化する。
各内部次数は有界有限自由$R_N$-複体であり、剰余は通常のKoszul双対性を満たす。
単元成分の消去によって
\[
 \mathbf R\underline{\operatorname{Hom}}_A(R_N,A)
 \simeq R_N(d)[-d]
\]
となる。中心正則元のchange-of-ringsから
\[
 \mathbf R\underline{\operatorname{Hom}}_B(R_N,B)
 \simeq R_N(a)[-(d-1)]
\]
を得る。この二つの計算は基底の有限global dimensionを用いない。

標準の導来section--torsion三角を用いる。
ここで必要な$\chi$はArtin基底でも確認できる。
$N>0$では$B/(h)$の$B$-自由分解は、中央の$h,h^N$による周期分解である。
有限生成$B$-加群$M$について、そのHom cohomologyは有限生成$B/(h)$-加群となる。
change-of-rings spectral sequenceと可換noetherian環$B/(h)$の$\chi$から
$\operatorname{Ext}_B^q(k,M)$は各$q$で内部次数について有限となる。
$R_N$の有限な$k$-filtrationにより、$\operatorname{Ext}_B^q(R_N,M)$も同じ性質を持つ。
$N=0$では通常の可換計算である。
従ってArtin--Zhangのsectionとlocal cohomologyの枠組を適用できる
\cite{LP-AZ}。

$M_\ell=B/B_{\ge\ell}$の内部次数filtrationの商は
$B_j\otimes_{R_N}R_N(-j)$である。
直前のrelative Gorenstein計算と$B_j$の自由性から
\[
 H^q(\mathbf R\tau B)_t=
 \begin{cases}
 B_{-t-a}^{\vee}&q=d-1,\\
 0&q\ne d-1
 \end{cases}
\]
を得る。固定した$t$では$j=-t-a$の一層のみが寄与し、$\ell>j$で安定する。
これは補題\ref{lp:lem:local-cohomology}と同じ有限次数の計算である。
三角$\mathbf R\tau B\to B\to\mathbf R\omega\pi B$から
\eqref{am:eq:line-hom}を得る。

自然な底変更写像についても確認する。
$M_\ell$の有限階数$B$-自由分解を、syzygyが各内部次数で$R_N$-自由となるように取る。
固定した内部次数でのHom複体は、各項が有限自由$R_N$-加群である。
その剰余はcohomologyが一つの次数にしかない複体であり、
体上のsplit contractionをnilpotent perturbationによって持ち上げると、
有限自由cohomologyの一項複体への実際のhomotopy同値を得る。
したがってこの複体はK-flatであり、通常のtensorが導来tensorを計算する。
項ごとのHomとtensorの同定、filtered colimit、上のsection三角を通じて
自然な底変更写像そのものが準同型であることが従う。
filtered colimitはK-flat複体をK-flat複体へ送るので、ここでunboundedな自由複体の
K-flat性を無条件に用いてはいない。

最後に$\mathbf R\omega\mathcal O(j)$の正のcohomologyは内部次数
$t\le-j-a$にしかない。
$i=1-a>-j-a$なので、$j\ge0$では内部切断後に全て消える。
次数零のsectionは$B(j)$であるから\eqref{am:eq:window-sections}を得る。
\end{proof}

\begin{lemma}[線束生成部分の有限模型]\label{am:lem:linear-category}
\[
 E=\bigoplus_{j=0}^{d-1}\mathcal O(j),\qquad
 \mathcal C_B=\operatorname{thick}_{R_N}\{\mathcal O(j):j\in\ZZ\}
\]
とすると$\mathcal C_B=\operatorname{thick}_{R_N}(E)$である。
全ての対象対のHomは$R_N$-perfectであり、Artin商への底変更と両立する。
例外成分$\langle\mathcal O,\ldots,\mathcal O(a-1)\rangle$を除く右直交を
\[
 \mathcal K_B^{\rm lin}
 =\mathcal C_B\cap\langle\mathcal O,\ldots,\mathcal O(a-1)\rangle^\perp
\]
とすると、有限mutationの射影$p$により$G=p(E)$がこれを生成する。
\end{lemma}

\begin{proof}
式\eqref{am:eq:relative-koszul}を$B$へ導来tensorする。
そのcohomologyは$R_N$と$R_N(-3)$のshiftだけであり、いずれも射影商で零となる。
従って線束の長さ$d$のKoszul型複体は射影商で完全である。
両端の係数の階数が一であるため、任意の連続した$d$個のtwistから全てのtwistを
前後に生成できる。
式\eqref{am:eq:line-hom}と有限cone・直和因子によりproper性と底変更が従う。
同式から$\mathcal O,\ldots,\mathcal O(a-1)$は$R_N$-例外列であり、
Homおよびその$R_N$-双対を用いて左右随伴が作れる。
従って有限個のevaluation coneが射影$p$を与え、$p(E)$が右直交を生成する。
\end{proof}

\subsection{MFのHomのproper性と生成子の持上げ}

\begin{lemma}[nilpotent perturbationによるMFの底変更]\label{am:lem:mf-hom}
有限自由graded MFの任意の二対象$P,Q$に対し、
$\operatorname{Hom}_{\MF}(P,Q)$は$R_N$-perfectである。
自然なHomの導来底変更は通常のMFの底変更Homと準同型である。
\end{lemma}

\begin{proof}
式\eqref{cg:eq:bidegree}のbigradingでHomを書けば、
各cohomological degreeは有限個の$A_j$の直和であり、有限自由$R_N$-加群である。
負のcohomological方向では必要な内部次数が負となるため零である。
剰余のHom複体$C_0$は、滑らかな古典cubicのgraded MFにおけるHomであり、
通常のOrlov同値により有限次元の有界cohomologyを持つ。

体上で$C_0$から有限cohomology複体$H$へのsplit contractionを選ぶ。
各次数の基底を持ち上げて、元のHom微分を$d_0+\delta$、$\delta\in h\operatorname{End}$と書く。
contractionのhomotopyを$s$とすると、
$(1-\delta s)^{-1}$および$(1-s\delta)^{-1}$は$h^{N+1}=0$により有限和である。
通常のperturbation恒等式をこの有限和に適用すると、
元のHom複体と$H\otimes_kR_N$上の変形微分との間のhomotopy同値を得る。
後者は有界有限自由複体である。
この同値からK-flat性が従い、項ごとのMF Homの底変更が導来底変更と一致する。
\end{proof}

\begin{lemma}[Artin基底上の生成子の持上げ]\label{am:lem:generator}
$P\in\MF^{\ZZ}(A,\Omega)$の剰余$\overline P$が古典graded MF圏を生成するとする。
このとき$P$は$\MF^{\ZZ}(A,\Omega)$を生成する。
\end{lemma}

\begin{proof}
補題\ref{am:lem:mf-hom}により、カテゴリーの底変更から古典MFへの関手は
Hom上で完全忠実である。像は$\overline P$を含むので、冪等完備化後にはMorita同値となる。
生成性を正確に持ち上げるため、$\mathcal D=\operatorname{Ind}(\MF^{\ZZ}(A,\Omega))$へ移る。
$P$はcompactであり、$\operatorname{RHom}(P,-)$はcolimitと可換である。
従って$M\in\mathcal D$が$P$と右直交なら
\[
 \operatorname{RHom}(\overline P,M\otimes_{R_N}^{\mathbf L}k)
 \simeq\operatorname{RHom}(P,M)\otimes_{R_N}^{\mathbf L}k=0.
\]
古典fiberでの生成性から$M\otimes_{R_N}^{\mathbf L}k=0$となる。
$R_N$はその$h$-filtrationによって$D(R_N)$内で$k$の有限extensionから作られる。
従って$R_N\in\operatorname{thick}(k)$であり、$M\otimes^{\mathbf L}k=0$は$M=0$を含意する。
よって$P$は$\mathcal D$のcompact generatorであり、compact部分を取って主張を得る。
ここではsmallなproper圏における右直交の消滅だけから生成性を結論していない。
\end{proof}

\subsection{有限windowによるenhanced比較}

\begin{theorem}[Artin線束残余pairのMF表示]\label{am:thm:comparison}
本節の入力に対し、有限window
\begin{equation}\label{am:eq:window-functor}
 \Psi=\operatorname{stab}\circ
       \bigl((\mathbf R\omega -)_{\ge\,4-d}\bigr)
\end{equation}
は$\mathcal K_B^{\rm lin}$上でdg Morita同値
\begin{equation}\label{am:eq:residual-mf}
 \Psi:\mathcal K_B^{\rm lin}\xrightarrow{\sim}\MF^{\ZZ}(A,\Omega)
\end{equation}
を与える。これはArtin商への底変更と両立する。
さらに$\rho=L_{\mathcal O}\circ(1)$に対し自然に
\begin{equation}\label{am:eq:rotation}
 \Psi\rho\simeq(1)\Psi
\end{equation}
である。特に$\rho$は自己同値となり、MFの三次周期性は
$\rho^3\simeq[2]$を与える。
\end{theorem}

\begin{proof}
まず\eqref{am:eq:window-functor}をdg quasi-functorとして定義する。
derived sectionと内部切断はenhanced関手であり、
stabilizationは補題\ref{am:lem:stabilization}で構成したdg Morita同値の逆である。
この関手が$\mathcal C_B$上で定義されることを生成子で確認する。
補題\ref{am:lem:line-bundles}から、$0\le j\le d-1$について
\[
 (\mathbf R\omega\mathcal O(j))_{\ge i}=B(j)_{\ge i}.
\]
これは$A$上perfectである。
実際、$B(j)$は$A$上長さ一の分解を持ち、有限長商
$T_j=B(j)_{<i}$の次数filtrationの各商は有限自由$R_N$のshiftである。
$R_N$は式\eqref{am:eq:relative-koszul}により$A$上perfectなので、$T_j$も同様である。
よって短完全列
\[
 0\to B(j)_{\ge i}\to B(j)\to T_j\to0
\]
の三項はいずれも$\mathcal D_{B/A}$に入る。
有限coneと直和因子へ延長して$\mathcal C_B$全体に$\Psi$を定義できる。

$0\le j<a$なら$-j\ge i$であり、$B(j)_{\ge i}=B(j)$は$B$上perfectである。
従って$\Psi$は例外成分$\mathcal O,\ldots,\mathcal O(a-1)$を零へ送る。
残る三つの$j=a,a+1,a+2$では
\begin{equation}\label{am:eq:three-images}
 \Psi(\mathcal O(j))\simeq\operatorname{stab}(T_j)[-1],\qquad
 \operatorname{gr}(T_j)=
 \bigoplus_{r=0}^{j-a}B_r\otimes_{R_N}R_N(j-r).
\end{equation}
従って三つの像は、三角的に
$\operatorname{stab}(R_N(a)),\operatorname{stab}(R_N(a+1)),
\operatorname{stab}(R_N(a+2))$を生成する。

剰余上では$\Psi$は通常のOrlovのsection--window関手である。
window $i=1-a$のkernelは$\langle\mathcal O,\ldots,\mathcal O(a-1)\rangle$であり、
その右直交への制限は古典Kuznetsov成分とgraded MFの同値である
\cite[Theorem 6.4 and its proof]{LP-BS}。
右直交と、同定されたwindowの反対側の直交の向きは混同しない。
両者は例外成分による同じVerdier商を表示し、$\Psi$の制限はその商関手である。

全ての自然なHom写像を剰余へ移すと、この古典同値のHom写像となる。
sourceのHomのproper性・底変更は補題\ref{am:lem:linear-category}、
targetは補題\ref{am:lem:mf-hom}による。
$\Psi$の底変更自然性はsectionの計算、内部切断、
補題\ref{am:lem:stabilization}の余核構成から従う。
従ってHom写像のconeは$R_N$-perfectで剰余が零であり、導来Nakayamaにより零である。
これで$\Psi|_{\mathcal K_B^{\rm lin}}$の完全忠実性を得る。
$G=p(E)$の像の剰余は古典MFのgeneratorである。
補題\ref{am:lem:generator}により$\Psi(G)$はArtin上でも全MFを生成するので、
本質的全射性が従う。

最後にrotationの自然性を確認する。
まず$\rho$は$\mathcal K_B^{\rm lin}$を保つ。
実際、$1\le j<a$では
$\operatorname{RHom}(\mathcal O(j),M(1))
=\operatorname{RHom}(\mathcal O(j-1),M)=0$であり、
式\eqref{am:eq:line-hom}より$\operatorname{RHom}(\mathcal O(j),\mathcal O)=0$である。
従ってevaluationの三角から$j\ge1$の直交性が従い、$j=0$の直交性はleft mutationの定義による。
$M\in\mathcal K_B^{\rm lin}$なら
\[
 (\mathbf R\omega M)_i
 \simeq\operatorname{RHom}(\mathcal O(-i),M)
 =\operatorname{RHom}(\mathcal O(a-1),M)=0.
\]
従って内部切断の自然な三角から
\[
 ((\mathbf R\omega M)(1))_{\ge i}
 =(\mathbf R\omega M)_{\ge i+1}(1)
 \simeq(\mathbf R\omega M)_{\ge i}(1)
\]
を得る。stabilizationは内部shiftを保つので、
$\Psi(M(1))\simeq\Psi(M)(1)$となる。
また$\Psi(\mathcal O)=0$であるから、evaluationの三角を$\Psi$で送ると
$\Psi(L_{\mathcal O}M(1))\simeq\Psi(M(1))$を得る。
これらの自然同型を合成して\eqref{am:eq:rotation}を得る。
ここでの周期同型は、指定された中心元のMFにおける周期変数$t$から移したものである。
\end{proof}

\subsection{正則局所模型からの底変更との関係}

\begin{proposition}[既存の$\Gamma_R$模型のArtin底変更]\label{am:prop:Gamma-basechange}
第\ref{lp:sec:main}節の正則局所基底$R$上の入力から、
剰余体が$\CC$である局所準同型$R\to R_N$によって$A,B$を底変更したとする。
このとき自然に
\[
 \operatorname{Perf}(\Gamma_R\otimes_R^{\mathbf L}R_N)
 \simeq\mathcal K_{B\otimes_RR_N}^{\rm lin}
 \simeq\MF^{\ZZ}(A\otimes_RR_N,\Omega\otimes1)
\]
であり、mutationで定めたrotation候補はMFの内部shiftへ移る。
\end{proposition}

\begin{proof}
既存の命題\ref{lp:prop:basechange}の結論を、変更先が非正則である場合へ無断で適用しない。
線束だけに必要な底変更を改めて確認する。
正則$R$上で$M_\ell=B/B_{\ge\ell}$の自由分解から作った、固定内部次数のHom複体は、
同命題の証明により有限自由cohomologyへhomotopy同値である。
従って$R_N$への通常のtensorが導来tensorを計算する。
syzygyの次数別$R$-分裂により、底変更後の自由分解も完全である。
そのHom、filtered colimit、section三角を用いると、
変更先について補題\ref{am:lem:line-bundles}で計算したHomと自然に一致する。
よって線束間のHomとその積について必要な底変更を得る。

$R$は正則なので、$-\otimes_R^{\mathbf L}R_N$は有界複体に対して有限Tor振幅を持ち、
射影商への関手が定義される。
線束間のHomの底変更を有限cone・直和因子・evaluationへ延長すれば、
$E_R,G_R$とそのendomorphism dg代数の底変更を得る。
これが第一の同値である。第二の同値とrotationの記述は
定理\ref{am:thm:comparison}による。
この証明は、$D^b(\operatorname{qgr}B)$全体の任意の底変更についての主張を含まない。
\end{proof}

\begin{remark}[結論の位置づけ]
第\ref{fo:sec:formal-cubic}節と本節を合成すると、上の入力を満たす全有限Artin切断において、
線束生成残余pairは同じ可換形式三次式$W_h$の残余pairへ移る。
非零の複素パラメータへの代数化、形式同値の収束、
および局所正則底上の全coherent残余圏と線束生成部分圏の一致は別問題である。
また本節のstabilizationは周囲でperfectな相対商を用いるので、
例\ref{cg:ex:artin-singularity}の絶対特異圏の反例と矛盾しない。
\end{remark}


\begin{corollary}[全有限Artin切断における可換残余復元]
\label{am:cor:formal-residual}
$d\ge5$とし、定理\ref{fo:thm:formal-reconstruction}の入力を取る。
$B_N=A_N/(\Omega_N)$、
$Y_N=V(W_h\bmod h^{N+1})\subset\PP^{d-1}_{R_N}$と置く。
このとき同じ形式三次式$W_h$を用いて、全ての$N$で
\[
 (\mathcal K_{B_N}^{\rm lin},\rho)
 \simeq
 (\Ku_{\rm perf}(Y_N/R_N),\rho_{Y_N})
\]
を得る。これは指定された内部shift、三次周期性、Artin商への底変更と両立する。
ここで$\Ku_{\rm perf}$は$\operatorname{Perf}(Y_N)$内で
$\mathcal O,\ldots,\mathcal O(d-4)$を除いた右直交である。
既存の正則局所模型からの入力には、命題\ref{am:prop:Gamma-basechange}も適用できる。
\end{corollary}
\begin{proof}
定理\ref{am:thm:comparison}を非可換側と可換側に適用し、
その間に定理\ref{fo:thm:formal-reconstruction}のMF同値を挟む。

可換側では、$\mathcal O,\mathcal O(-1),\ldots,\mathcal O(-(d-1))$の直和が
$D_{\rm QCoh}(Y_N)$をcompactに生成する
\cite[Tag 0BQT, Remark 36.16.4]{AMStacks}。
具体的には、周囲の射影空間の相対対角分解がその線束生成子を与え、
閉埋込みのpushforwardがconservativeであることと随伴により、
そのpullbackが$Y_N$の生成子となる。
qcqs schemeでcompact対象はperfect対象である
\cite[Tag 09M1, Proposition 36.17.1]{AMStacks}から、
これらの線束の厚い閉包は$\operatorname{Perf}(Y_N)$である。
補題\ref{am:lem:linear-category}で用いた全twistの厚い閉包と一致する。
各関手のrotationと底変更の両立から、主張するpairの同値を得る。
\end{proof}
